\documentclass[11pt]{article}
\usepackage[a4paper,margin=25mm]{geometry}
\usepackage[T1]{fontenc}
\usepackage{lmodern,microtype,amsmath,amssymb,amsthm,mathtools}
\usepackage[colorlinks=true,linkcolor=blue,urlcolor=blue,citecolor=blue]{hyperref}
\newtheorem{proposition}{Proposition}
\newtheorem{remark}[proposition]{Remark}
\newcommand{\Ric}{\operatorname{Ric}}
\newcommand{\CP}{\mathbb{CP}}
\newcommand{\calC}{\mathcal C}
\title{Q3: The Zero Mode and the Proposed Horikawa Candidate}
\author{Continuation of the supplied October 8 research}
\date{8 October 2026}
\begin{document}
\maketitle
\begin{abstract}
We continue the search for a smooth closed non-Bach-flat absolute minimizer
of $P(g)=\int|\Ric_g^\circ|^2$. The newly proposed Horikawa surface gives
a valid fixed attainment problem, but its defining equation explicitly
contains both $(-3)$-sphere and ADE vanishing-cycle topology. For the
surviving $\mathcal O(-3)$ degeneration of the previous continuation,
we solve the zero-mode radial inner problem exactly, including its finite
renormalized quadratic minimum. We also compute the mixed bulk variation:
an interior Euler--Lagrange term cancels an apparent point-jet shortcut.
The inner minimum is negative, but the sign of the complete closed-manifold
replacement is still undetermined. No smooth attainment theorem is asserted.
\end{abstract}

\section{The fixed target and the new candidate}
The target is a smooth absolute minimizer of $P$ which is not Bach-flat.
It suffices to prove attainment on a fixed manifold where all conformal
classes are Yamabe-negative and no Einstein metric exists. Indeed a closed
$P$-critical metric has constant scalar curvature; when this constant is
nonzero, Bach-flatness and $P$-criticality together imply Einstein.
These variational facts are inherited from the preceding work.

The user-provided proposal is
\[
 X_0=\{w^2+z^5+x^{10}+y^{10}=0\}\subset\mathbb P(1,1,2,5),
 \qquad M_H=X_0\#\overline{\CP^2}.
\]
The affine-cone gradient vanishes only at the origin; $X_0$ avoids the
two singular coordinate points of the ambient weighted projective space.
Weighted adjunction gives $K=\mathcal O_{X_0}(1)$ and
$K^2=10/(1\cdot1\cdot2\cdot5)=1$. The standard weighted hypersurface
calculation gives $p_g=2,q=0$; this is the classical $(1,2)$ family
\cite{CP}. Consequently
\[
 \chi(M_H)=36,\qquad \tau(M_H)=-24,\qquad
 (b^+,b^-)=(5,29),\qquad 2\chi+3\tau=0.
\]
The ordinary sign classes $K\pm E$ provide a projection budget at least
one for every metric. Combining the classical mixed monopole curvature
bound \cite{LeBrun} with the supplied strict negative-KE insertion gives
\begin{equation}
 \frac{16\pi^2}{3}\le\inf_{g\in\operatorname{Met}(M_H)}P(g)<8\pi^2.
 \label{eq:interval}
\end{equation}
The upper bound here is an inherited project result, not re-proved in
this note. The positive actual monopole budget excludes nonnegative
scalar-curvature representatives. Thus every conformal class is
Yamabe-negative, and the lower bound excludes Einstein metrics.
Horikawa therefore gives the same attainment reduction as the previously
selected Godeaux blowup, with the advantage $b^+>1$ for ordinary
Seiberg--Witten chamber bookkeeping.

\begin{proposition}[Explicit topology in this particular Horikawa model]
The surface $X_0$ contains an embedded smooth $(-3)$-sphere $C$ with
$K\cdot C=1$. It also contains five disjoint $A_4$ Milnor-fibre regions;
their vanishing-cycle lattices pair trivially with $K$. The blowup can
be taken away from these configurations.
\end{proposition}
\begin{proof}
Choose $c^5=-1$ and $\epsilon\in\{1,-1\}$. The map
\[
 [u:v]\longmapsto[u:v:cv^2:\epsilon i u^5]
\]
respects the coordinate weights and lies in $X_0$, since
$-u^{10}-v^{10}+u^{10}+v^{10}=0$. Its first two coordinates recover
$[u:v]$, so it embeds $\CP^1$. The curve lies in the weight-one charts,
and $\mathcal O(1)$ restricts with degree one. Adjunction gives
$-2=C^2+K\cdot C$, hence $C^2=-3$.

For the second assertion, consider the proper family
\[
 X_a=\{w^2+z^5+x^{10}+y^{10}+ax^5y^5=0\}.
\]
The singularity equations force $w=z=0$, $x,y\ne0$, and
$2x^5+ay^5=2y^5+ax^5=0$. Thus the only singular parameters are
$a=\pm2$. At $a=2$ there are five singular points,
$w=z=0$, $x^5+y^5=0$. In local analytic coordinates their equation is
\[
 u^2+v^2+z^5=0,
\]
the $A_4$ singularity. A nearby smooth fibre contains five disjoint
Milnor fibres with the negative $A_4$ intersection matrix. Smooth fibres
are diffeomorphic along paths in $\mathbb C\setminus\{-2,2\}$, giving
these regions in $X_0$. Their local canonical bundle is trivial, so $K$
pairs to zero on their vanishing spheres. A blowup outside the chosen
regions leaves all these pairings unchanged.
\end{proof}

The $(-3)$ filling supports scalar-flat K\"ahler ALE metrics with
$P=8\pi^2/3$, while the $A_4$ topology supports hyperk\"ahler ALE
metrics with $P=0$ \cite{Kronheimer,Suvaina}. This proves that these
topological possibilities are present; it does not prove that minimizing
sequences develop them. Neither the interval \eqref{eq:interval} nor the
displayed sign-class pairings alone excludes them. A change to Horikawa
therefore does not supply the missing compactness theorem.

\section{The remaining first-order zero mode}
The previous continuation combined classical ALE mass flexibility with
the supplied A02.59D boundary flux. For a marked $\mathcal O(-3)$
replacement at a smooth $\frac13(1,1)$ orbifold point of a $P$-critical
bulk, it obtained
\begin{equation}
 P(g_{\lambda,m})=P(g)+\frac{8\pi^2}{3}
       -24\pi^2m\mu\lambda^2+o(\lambda^2),
 \label{eq:mass}
\end{equation}
where $m$ is Hein--LeBrun mass and
$\mu=\langle(W^++sI/12)\zeta,\zeta\rangle$ on the invariant
self-dual axis. Both mass signs are available on the same marked filling
at the same complete $P$-energy. Under the stated marked local convergence
and replacement hypotheses, minimality thus requires $\mu=0$.

The scalar isotropy then gives
\[
 R^+=W^++\frac{s}{12}I=\operatorname{diag}(\nu,\nu,0),
 \qquad \nu=s/8.
\]
The previously computed radial logarithmic term vanishes at this block.
The next calculation treats its finite inner term. It uses the radial
LeBrun filling, not the nonradial Stein fillings used for mass flexibility.

\section{Exact solution of the radial inner problem}
Normalize the $\mathcal O(-3)$ bolt radius to one and write
\[
 g_0=\frac{dr^2}{G}+r^2(\sigma_1^2+\sigma_2^2+G\sigma_3^2),
 \qquad G=1+r^{-2}-2r^{-4}.
\]
The angular volume is $2\pi^2/3$. For
\[
 F=G+tu+O(t^2),\qquad H=G+tv+O(t^2),\qquad d=u-v,
\]
use $g_t=dr^2/H+r^2(\sigma_1^2+\sigma_2^2+F\sigma_3^2)$.
The previously derived radial Cartan formulas give the linearized
self-dual curvature eigenvalues
\begin{align}
 L_1&=\frac{d'}{4r}+\frac d{2r^2},\label{eq:L1}\\
 L_3&=-\frac{u''}{4}-\frac{7u'}{4r}-\frac{2u}{r^2}+L_1
 +\frac{d'G'}{8G}
 +d\left(\frac{G'}{rG}+\frac{G''}{4G}
                    -\frac{(G')^2}{8G^2}\right).\label{eq:L3}
\end{align}
Smoothness at the fixed bolt requires $F(1)=H(1)=0$ and
$F'(1)H'(1)=36$. At first order this reads
\[
 u(1)=v(1)=0,\qquad u'(1)+v'(1)=0.
\]
Higher-order corrections can enforce the exact bolt product; they do not
change the quadratic chiral energy on any fixed compact inner region,
since $R^+(g_0)=0$.

\begin{proposition}[Radial linear response]
Let $x=r^2$. The smooth-bolt linear problem $L_1=\nu$, $L_3=0$
has the solution
\begin{align}
 d_0&=\nu(x-x^{-1}),\label{eq:d}\\
 u_0&=\nu\left[
 \frac x6-\frac12+
 \frac{2+\tfrac12\log((x+2)/3)}x+
 \frac{-\tfrac53-2\log((x+2)/3)}{x^2}\right],\label{eq:u}\\
 v_0&=u_0-d_0.\nonumber
\end{align}
It has the prescribed quadratic background
$u_0\sim\nu x/6$, $v_0\sim-5\nu x/6$.
\end{proposition}
\begin{proof}
Equation \eqref{eq:L1} is $(d_x+d/x)/2=\nu$. The bolt value fixes
the integration constant and gives \eqref{eq:d}. Substitution in
\eqref{eq:L3} reduces $L_3=0$ to
\[
 xu_{xx}+4u_x+2u/x
   =\nu\frac{2x^2+6x+1}{2(x+2)^2},
 \qquad u(1)=0,\quad u_x(1)=\nu.
\]
The substitution $u=y/x$ gives an elementary second-order equation;
integration yields \eqref{eq:u}. Direct differentiation verifies
$u_x(1)=\nu$, $v_x(1)=-\nu$ and the equation. At infinity,
\[
 u_0=\nu\left[\frac x6-\frac12+
 \frac{\frac12\log x+2-\frac12\log3}{x}
 +\frac{-2\log x-\frac23+2\log3}{x^2}+O(x^{-3})\right].
\]
In particular the constant first correctors agree, as required by the
vanishing logarithmic chiral-energy term.
\end{proof}

\section{The exact finite inner minimum}
Use the operator norm for Weyl curvature and put
\[
 \calC=\int(4|W^+|^2+s^2/12)\,dV=4\int|R^+|^2\,dV.
\]
On a closed manifold $P=\calC-8\pi^2(2\chi+3\tau)$.
For $X=R^2$, the coefficient of $t^2$ in the inner chiral integral is
\[
 \mathcal Q_X(u,d)=\frac{2\pi^2}{3}
 \int_1^X\left[x(d_x+d/x)^2+2xL_3^2\right]dx.
\]
The corresponding background orbifold ball contributes
$(2\pi^2/3)2\nu^2X^2$. Define the renormalized coefficient by
subtracting this expression and then letting $X\to\infty$.

\begin{proposition}[A minimum with fixed full asymptotic data]
For smooth-bolt competitors for which
$z=x(d-d_0)$ satisfies $z(1)=z(\infty)=0$ and the following integrals
converge,
\begin{equation}
 \mathcal Q_{\rm ren}(u,d)+\frac{4\pi^2}{3}\nu^2
 =\frac{2\pi^2}{3}\int_1^\infty
       \left[\frac{(z')^2}{x}+2xL_3(u,d)^2\right]dx\ge0.
 \label{eq:square}
\end{equation}
Equality holds for \eqref{eq:d}--\eqref{eq:u}. Thus the exact minimum is
\begin{equation}
 \boxed{\mathcal Q_{\rm ren,min}=-\frac{4\pi^2}{3}\nu^2
                              =-\frac{\pi^2s^2}{48}.}
 \label{eq:inner}
\end{equation}
\end{proposition}
\begin{proof}
Since $d_x+d/x=2\nu+z'/x$,
\[
 x(d_x+d/x)^2=4\nu^2x+4\nu z'+(z')^2/x.
\]
Integration cancels the cross term by the two endpoint values of $z$.
For the displayed solution, $L_3=0$ and
$\mathcal Q_X=(2\pi^2/3)2\nu^2(X^2-1)$, proving the finite constant.
The remaining homogeneous solutions for $u$ are $x^{-1}$ and $x^{-2}$;
the zero bolt value and required derivative fix their coefficients.
\end{proof}

\paragraph{Boundary dependence.}
If $d=\nu x+c/x+o(x^{-1})$ instead, the right side of
\eqref{eq:square} acquires
\[
 \frac{2\pi^2}{3}\,4\nu(c+\nu).
\]
Thus the matching data at infinity are essential to the finite minimum.
Equation \eqref{eq:inner} is an inner, successive-limit calculation:
first differentiate at $t=0$ on each fixed inner ball, then take its
radius to infinity. For squared physical scale $\varepsilon=\lambda^2$
it corresponds to an inner contribution
$-\pi^2s^2\varepsilon^2/48$. It is not yet the coefficient of a
complete closed-manifold competitor. The outer response and the
subleading matching terms have the same order. For example, simply
cutting off a mass tail at radius $r_0$ allows an error of size
$O(\varepsilon^2r_0^{-4})$. The previously proved $o(\lambda^2)$
remainder in \eqref{eq:mass} does not control order $\lambda^4$.

\section{The full mixed variation and its cancellation}
One might try to change the bulk curvature jet slightly so that
$\mu\ne0$ and then apply \eqref{eq:mass}. This loses bulk criticality.
The missing term can be written explicitly.

Let $E_g$ denote the variational gradient of $\calC$, with
$D\calC_g[h]=\int\langle E_g,h\rangle$. Use the invariant mass tensor
\[
 B(x)=\frac{r^2I-xx^{\mathsf T}-(J_0x)(J_0x)^{\mathsf T}}{r^4},
 \qquad H_m=9m\chi B,
\]
with a fixed cutoff $\chi=1$ near the puncture. The punctured first
variation which gave \eqref{eq:mass} has, for a noncritical bulk,
\begin{equation}
 a(g,m)=-24\pi^2m\mu(g,p)+\int\langle E_g,H_m\rangle\,dV.
 \label{eq:noncritical}
\end{equation}
The latter integral is locally integrable. At a critical base metric,
writing $L_g=DE_g$, its derivative is
\begin{equation}
 \ell_m(h)=D_ga[h]=-24\pi^2mD\mu_g[h]
                         +\int\langle L_gh,H_m\rangle\,dV.
 \label{eq:mixed}
\end{equation}
Contractions, volume and transported coordinate choices contribute no
additional gradient term there, because $E_g=0$.

Here is the consequence for a concentrated jet. Let
$h_r(x)=r^2h_1(x/r)$, with $h_1$ invariant, compactly supported and
quadratic near zero. Its curvature variation $D\mu[h_r]$ can stay
nonzero as $r\downarrow0$. If $K_0=D(R^+)_\delta$, then
\[
 D^2\calC_g[h_r,h_r]
       =8r^4\int|K_0h_1|^2+O(r^6).
\]
Nevertheless the complete coefficient satisfies
\begin{equation}
 |\ell_m(h_r)|\le C|m|r^2.\label{eq:localbound}
\end{equation}
Indeed $K_0B=0$ away from zero: $B$ is the infinitesimal scalar-flat
K\"ahler tail. In normal coordinates $g-\delta=O(|x|^2)$, so
$D(R^+)_g[H_m]=O(|m||x|^{-2})$. Its product with $K_gh_r$ integrates
to $O(|m|r^2)$. The other mixed terms are a bounded background
curvature times $h_r\partial^2H_m$, $H_m\partial^2h_r$, or
$\partial h_r\partial H_m$; they have the same integrable bound.
This accounts for the moving norms and volume as well. Equivalently,
integration by parts in the second term of \eqref{eq:mixed} cancels
the apparent point-evaluation term in the first.

For a nonzero principal quadratic form the resulting ratio is bounded:
\[
 \frac{|\ell_m(h_r)|^2}{D^2\calC_g[h_r,h_r]}=O(m^2).
\]
It does not diverge like $r^{-4}$. Thus a local point-jet variation
cannot be used to overwhelm every finite gluing coefficient.

\section{The actual remaining comparison}
If a complete uniform expansion can be established at $\mu=0$,
it must have the form
\[
 P(G_{\lambda,t})-P(g)-\frac{8\pi^2}{3}
 =\frac12t^2Q(h,h)+t\lambda^2\ell_m(h)+\lambda^4b_m
           +o(t^2+|t|\lambda^2+\lambda^4).
\]
For $Q(h,h)>0$, optimizing $t$ then gives the finite coefficient
\[
 b_m-\frac{\ell_m(h)^2}{2Q(h,h)}.
\]
The negative number in \eqref{eq:inner} supplies one computable inner
part, under its specified matching conditions. It does not determine
$b_m$, which must include the exterior Jacobi response, all boundary
terms and the higher background jets. A nonnegative Hessian alone does
not determine this total sign. No uniform fourth-order expansion or
negative total coefficient has been established here.

For completeness, a return to purely explicit K\"ahler or homogeneous
ansatzes does not give an immediate alternative. A closed K\"ahler
$P$-critical metric has constant scalar curvature. When $s\ne0$, varying
the K\"ahler class makes $(c_1\cdot\Omega)^2/\Omega^2$ stationary,
forcing $c_1\parallel\Omega$; harmonicity of the Ricci form then makes
the metric Einstein. When $s=0$ it is scalar-flat ASD. The complete
homogeneous classification \cite{BCGV}, Remark 1.5, likewise implies
that a locally homogeneous $P=F_{-1/4}$ critical four-metric is Einstein
or an opposite-curvature surface product. These candidates are all
Bach-flat.

The fixed goal remains smooth global attainment. This continuation
establishes explicit topology in the proposed alternative surface and
an exact radial inner minimum; it also identifies the necessary mixed
variation. The zero mode, other ALE topologies and general degeneration
have not been excluded. No answer to Q3 is claimed.

\begin{thebibliography}{9}
\bibitem{CP} S. Coughlan and R. Pignatelli,
\emph{Simple fibrations in $(1,2)$-surfaces},
Forum Math. Sigma 11 (2023), e43.
\href{https://doi.org/10.1017/fms.2023.41}{doi:10.1017/fms.2023.41}.
\bibitem{LeBrun} C. LeBrun,
\emph{Ricci Curvature, Minimal Volumes, and Seiberg--Witten Theory},
\href{https://arxiv.org/abs/math/0003068}{arXiv:math/0003068}.
\bibitem{Kronheimer} P. B. Kronheimer,
\emph{The construction of ALE spaces as hyper-K\"ahler quotients},
J. Differential Geom. 29 (1989), 665--683.
\bibitem{Suvaina} I. \c{S}uvaina,
\emph{ALE Ricci-flat K\"ahler metrics and deformations of quotient
surface singularities},
\href{https://arxiv.org/abs/1102.2672}{arXiv:1102.2672}.
\bibitem{BCGV} M. Brozos-V\'azquez, S. Caeiro-Oliveira,
E. Garc\'ia-R\'io and R. V\'azquez-Lorenzo,
\emph{Four-dimensional homogeneous critical metrics for quadratic
curvature functionals},
\href{https://arxiv.org/abs/2309.02108}{arXiv:2309.02108}.
\bibitem{Prior} \emph{Q3: Mass Variation and a Cyclic Degeneration
Obstruction}, preceding October 8 continuation. Its first-order formula
uses the supplied A02.59D boundary-flux calculation. The radial operators
in \eqref{eq:L1}--\eqref{eq:L3} were derived in the preceding radial
LeBrun analysis. This note independently differentiates the displayed
solution and verifies the square-difference identity exactly.
\end{thebibliography}
\end{document}
